% Revised 2026-08-28 after full-text verification of farea2025qcpinn and
% hqpinn2025flow (see references.bib changelog comment). Numbers and quotes
% below are checked against the source PDFs, not abstract summaries.
% Remaining citations (qcpinn2025reservoir, hqcpinn2026hydro, qpinnmac2025)
% are still unverified stubs -- see references.bib TODOs.

Physics-informed neural networks (PINNs) \citep{raissi2019pinn} have become
a standard tool for solving forward and inverse PDE problems by embedding
the governing equations directly into the training loss. Motivated by
broader interest in quantum machine learning for scientific computing, a
growing line of work replaces part of the classical network with a
variational quantum circuit (VQC), typically in a hybrid architecture where
classical layers handle pre- and post-processing while the VQC provides the
trainable core -- an approach variously termed QPINN or QCPINN.

\citet{farea2025qcpinn} introduce QCPINN and evaluate it, alongside a
continuous-variable (qumode-based) variant that trains very poorly across
the board, on five forward-only PDE benchmarks with Dirichlet-type boundary
conditions: Helmholtz, a 2D lid-driven cavity flow, a 1D wave equation,
Klein-Gordon, and convection-diffusion. Their headline claim -- 10--30\% of
the trainable parameters of a classical PINN with comparable or better
accuracy -- holds for three of the five: a 4--64\% reduction in relative
$L_2$ error on Helmholtz, Klein-Gordon, and convection-diffusion. On the
remaining two, the quantum branch is \emph{less} accurate than the
classical baseline despite using far fewer parameters (cavity: 11--23\%
higher relative error across the three output fields; wave: 12\% higher),
which the authors state plainly rather than omit \citep[Sec.~8,
``Conclusion'']{farea2025qcpinn}. They also report a timing column absent
from the abstract: per-iteration wall-clock time is 15--40$\times$ that of
the classical baseline in every one of the five benchmarks, which they
attribute to the cost of simulating the circuit classically rather than to
the algorithm itself. Separately, their qubit-based circuit sweep tests
four named topologies (Alternate, Cascade, Cross-mesh, Layered) crossed
with two data-encoding schemes (amplitude, angle); the accuracy and timing
numbers above are for each benchmark's best-performing topology, most
often Cascade or Cross-mesh -- both of which use continuously-tunable
controlled-rotation entanglers (CRX/CRZ) rather than fixed CNOT gates,
which the authors argue is why those two topologies converge more reliably
than Alternate or Layered.

\citet{hqpinn2025flow} report a closely related pattern on a different
problem family: the compressible Euler equations, tested across a smooth
harmonic traveling wave, a non-harmonic solution with a moving contact
discontinuity, and a 2D transonic aerofoil flow with a shock. A quantum
branch whose functional form is an explicit truncated Fourier series
achieves the single best loss of any model on the harmonic case (an
order-of-magnitude improvement over the best classical network, using far
fewer parameters), but ``significantly underperformed'' on the non-harmonic
case and underperformed further still on the shock-dominated transonic
case, where the quantum branch is additionally under-parameterized. The
paper's own stated thesis is that harmonicity of the target solution, not
problem difficulty in general, is what determines whether the quantum
branch helps: ``quantum PINNs excel in harmonic regimes... due to the
Fourier structure of PQCs... [but] struggle to generalize to
discontinuities or shocks.'' On timing, wall-clock cost is reported only
for the transonic case, where per-epoch cost for the quantum configurations
ranges 17--82$\times$ that of the classical configurations tested (the
paper's own summary rounds this to ``up to 100$\times$''); a separate,
larger ``two orders of magnitude'' figure appearing earlier in the paper is
attributed to other authors' measurements, not to this paper's own
experiments. Notably, this paper does not use the term ``barren plateau''
anywhere; its account of trainability loss in larger quantum models is
framed instead as ``excessive expressibility,'' a distinct (if related)
diagnosis.

Read together, these two papers already show that the value of a quantum
branch in a hybrid PINN is problem-dependent rather than uniform -- within
a single paper's own benchmark suite, in both cases -- and that reporting
practices vary: parameter-count efficiency is reported everywhere, but
wall-clock or per-iteration time is reported only for a subset of
benchmarks in each paper (all five in \citealp{farea2025qcpinn}, only the
hardest of three in \citealp{hqpinn2025flow}), which limits how directly
either paper's timing claims generalize.
\todo{Add hydrological \citep{hqcpinn2026hydro}, reservoir seepage
\citep{qcpinn2025reservoir}, and the theoretical barren-plateau bound
\citep{qpinnmac2025} once those are verified against full text -- see
references.bib TODOs. Until then, do not restate the epoch-count /
polynomial-gradient-bound claims for these three beyond what the earlier
abstract-level search already summarized in chat, since that summary has
not been checked against the full papers the way the two above have.}

Our work differs from this pair of closely related studies in three ways.
First, both restrict the quantum branch to forward, Dirichlet-type problems
(\citealp{farea2025qcpinn} exclusively; \citealp{hqpinn2025flow} uses
Dirichlet conditions throughout as well); we test both forward and inverse
formulations across three boundary-condition types (Dirichlet, Neumann,
Robin) on a single diffusion PDE with closed-form solutions in all six
cases, which lets us measure error against ground truth exactly rather than
against a numerical or experimental reference. Second, where
\citet{farea2025qcpinn} report per-iteration time as one column among many
and \citet{hqpinn2025flow} report it for only one of three benchmarks, we
report wall-clock training time as a primary outcome for every scenario
alongside accuracy, and find the quantum branch uniformly worse on both
axes in all six -- a stronger and more consistent negative result than
either prior paper's own (partially mixed) findings. Third, neither prior
paper isolates circuit structure as an experimental variable with direct
gradient measurement: \citet{farea2025qcpinn} discusses trainability
differences between entangler choices (CRX/CRZ vs.\ fixed CNOT)
qualitatively, and formally invokes barren plateaus only for its
continuous-variable circuit, which is not the qubit-based design we build
on; \citet{hqpinn2025flow} attributes reduced trainability in larger
circuits to ``excessive expressibility'' based on final-loss curves, not a
direct gradient-norm measurement. We instrument the circuit-only gradient
norm directly across a controlled qubit-count/depth sweep, providing an
empirical data point neither paper's diagnostic approach supplies.
